% =============================================================================
%  GENERAL INTRODUCTION
%  File: chapters/General Introduction.tex
%  Do NOT add \documentclass — this file is \input{} from main.tex
% =============================================================================

\section*{General Context and Motivation}
\addcontentsline{toc}{section}{General Context and Motivation}

The contemporary retail enterprise operates at the intersection of two transformative forces: the exponential proliferation of consumer transaction data and the maturation of artificial intelligence techniques capable of extracting commercial intelligence from that data at scale. In this environment, the competitive boundary between firms is increasingly drawn not by product quality or store location alone, but by the analytical sophistication with which accumulated purchase histories are converted into forward-looking customer knowledge.

In the Algerian retail market, this transition is at an early but accelerating stage. The emergence of exclusive brand distributors — firms that manage the entire import, distribution, and retailing of a single international brand within the national territory — creates an analytically privileged observatory: a clean, single-brand transactional database free of the multi-brand confounds that complicate segmentation in hypermarket or marketplace environments. SARL Apparel, the exclusive distributor of PUMA in Algeria since 2018, exemplifies this structure. Its proprietary ERP dataset — over 180,000 transaction lines spanning multiple fiscal years, 35,696 identifiable customers, and three product lines (Footwear, Textile, Accessories) — constitutes the empirical foundation of this thesis.

\section*{Research Problem}
\addcontentsline{toc}{section}{Research Problem}

Despite the data richness of its operational environment, SARL Apparel's commercial decisions at the time of this study remained predominantly intuition-driven: customer contact campaigns were undifferentiated mass broadcasts, inventory replenishment relied on managerial experience rather than algorithmic demand signals, and no systematic estimation of individual customer lifetime value informed acquisition or retention budget allocation. This gap between data availability and analytical exploitation defines the central problem of this thesis:

\begin{quote}
\textit{Can an integrated, AI-driven customer intelligence pipeline — combining machine-learning customer segmentation, probabilistic Customer Lifetime Value modelling, and statistical demand forecasting — provide SARL Apparel with actionable, quantitatively validated instruments for transforming raw transactional data into sustainable competitive advantage?}
\end{quote}

\section*{Research Objectives and Hypotheses}
\addcontentsline{toc}{section}{Research Objectives and Hypotheses}

This thesis pursues three structured research objectives, each associated with a testable hypothesis:

\begin{enumerate}

    \item \textbf{Objective 1 — Segmentation}: Demonstrate that K-Means clustering applied to log-normalised RFM features produces a customer partition with meaningfully higher cluster validity (silhouette score, Davies-Bouldin index) than conventional quintile-based RFM scoring.

    \textbf{H\textsubscript{1}}: K-Means on RFM achieves a mean silhouette score $> 0.35$ and a DB index $< 1.5$ on the PUMA Algeria dataset, outperforming quintile RFM on both indices.

    \item \textbf{Objective 2 — CLV Modelling}: Validate the BG/NBD + Gamma-Gamma probabilistic CLV pipeline on the mixed B2B/B2C non-contractual transaction dataset and estimate a customer-level CLV distribution with managerial utility.

    \textbf{H\textsubscript{2}}: The BG/NBD model achieves a 90-day holdout MAE $< 0.5$ repeat transactions per customer on the PUMA Algeria panel.

    \item \textbf{Objective 3 — Forecasting}: Determine the optimal demand forecasting model for PUMA Algeria daily revenue from a set of four candidates (SARIMA, Prophet, XGBoost, Random Forest) under the specific conditions of (i)~a short observation history, (ii)~a strong weekly seasonal cycle, and (iii)~high intra-week amplitude oscillation.

    \textbf{H\textsubscript{3}}: SARIMA, by virtue of its explicit seasonal memory at period $m=7$, outperforms ML-based alternatives on the 60-day out-of-sample MAE.

\end{enumerate}

\section*{Thesis Structure}
\addcontentsline{toc}{section}{Thesis Structure}

The thesis is organised into three chapters, each corresponding to one analytical layer of the integrated pipeline:

\begin{itemize}[noitemsep]
    \item \textbf{Chapter 1: Theoretical Foundations} establishes the conceptual and mathematical framework — AI and machine learning taxonomy, RFM segmentation theory, and the probabilistic CLV pipeline (BG/NBD + Gamma-Gamma) — providing the formal vocabulary for the literature review and empirical chapters.

    \item \textbf{Chapter 2: Literature Review} provides a thematically organised synthesis of the empirical research on customer segmentation, CLV modelling, and demand forecasting. It identifies three specific gaps in the existing literature that the thesis addresses: the absence of RFM-K-Means benchmarks in North African exclusive-distribution retail contexts, the untested applicability of the BG/NBD pipeline to mixed B2B/B2C channel structures, and the lack of CLV-infused demand forecasting experiments in Algerian fashion retail.

    \item \textbf{Chapter 3: Methodology, Empirical Application, and Results} presents the complete data engineering pipeline, segmentation and CLV results, multi-model forecasting experiment, and a set of strategically actionable managerial recommendations grounded in the quantitative findings.
\end{itemize}